\begin{figure}[ht]
\centering
\begin{tikzpicture}[x=.82cm,y=.68cm,>=Stealth,font=\small]
 \fill[blue!6] (3.65,-.45) rectangle (4.35,3.5);
 \foreach \j in {0,...,8}{
  \node at (\j,3.8){\(\j\)};
  \foreach \h in {0,1.5,3}{\fill[black!35] (\j,\h) circle(1.1pt);}
 }
 \draw[->,very thick] (3,3)--(2,3);
 \draw[->,very thick] (2,3)--(0,3);
 \node[left] at (-.4,3){left};
 \draw[->,very thick,blue!65!black] (6,1.5)--(4,1.5);
 \draw[->,very thick,blue!65!black] (4,1.5)--(2,1.5);
 \draw[->,thick,dashed,red!65!black] (5,1.7) to[bend right=24] (3,1.7);
 \node[left] at (-.4,1.5){central};
 \draw[->,very thick] (8,0)--(7,0);
 \draw[->,very thick] (7,0)--(5,0);
 \node[left] at (-.4,0){right};
 \node[below,align=center] at (4,-.6){middle row \(m=4\)\\
           unused factor \(u\) on the dashed jump};
\end{tikzpicture}
\caption{The partition at the illustrative odd level \(N=9\).
A central path may visit row \(m\), or skip it by \(m+1\to m-1\).
The skipped row contributes its unused factor \(u\), retained in the
exact central polynomial even though it vanishes at \(u=0\).}
\label{lc:path-figure}
\end{figure}
